\chapter{总结与展望}

\section{工作总结}

到论文撰写完成时，Mini-Chat-Bar 已经完成了从需求分析、系统设计到实现与测试的主要工作。这个系统面向的不是泛社交聊天，而是技术交流场景里更常见的几条链路：用户登录后可以进入主聊天界面，在私聊和聊天室之间切换；消息支持文本、代码、图片和文件；需要反复回看的内容可以加入收藏；讨论过程中如果需要补充说明或快速整理结论，还可以调用 AI 问答和讨论总结功能。

结合整个毕业设计的推进过程，我完成的工作主要集中在几个连续阶段。前期先梳理需求、确定论文结构，并在此基础上把系统定为前后端分离架构；随后围绕用户、私聊、聊天室、收藏、AI 对话与聊天总结等模块整理业务流程和数据结构；在实现阶段继续完成 Web 端界面、Electron 桌面端适配、后端接口、Socket.IO 实时链路以及 AI 问答与总结功能的联调；最后再把需求分析、系统设计、实现过程和测试结果整理成论文。对我来说，这次毕设不是只把一个功能做出来，而是把分析、设计、编码、测试和文档整理完整地走了一遍。

从当前实现结果看，系统的主要链路已经可以闭环运行。前端已经完成登录、主聊天界面、聊天室详情、收藏页和 AI 面板等核心页面；后端已经完成用户认证、消息处理、聊天室管理、收藏管理和 AI 流式问答等服务；MongoDB 用于保存用户、房间、私聊消息、聊天室消息、收藏和 AI 相关数据。具体到业务流程上，邮箱登录后可以进入 \texttt{ChatView.vue}，私聊通过 \texttt{/api/chat/messages/:id} 与 \texttt{private-message} 联动，聊天室消息通过 \texttt{/room/:roomId/messages} 和 \texttt{group-message} 同步，\texttt{@AI} 提问则通过 \texttt{/api/chatroom-ai/ask-stream} 与 \texttt{ai-stream-chunk} 回写结果。

第六章的测试结果也说明，这套系统已经具备基本的演示和使用条件。普通接口保持毫秒级响应，私聊消息和聊天室消息可以稳定写入并同步，桌面端的托盘驻留、未读提醒和消息预览已经能够正常工作。AI 问答与总结链路虽然明显慢于普通接口，但已经跑通了从提问、检索、生成到结果落库的完整流程。对本科毕业设计而言，系统已经完成了预期的主要目标。

与本课题实施过程相关的任务书、开题/中期检查、指导记录、审批表和成绩评定等材料，已统一整理于\hyperref[app:materials]{附录}。

\section{实践体会}

这次毕业设计让我更直接地体会到，聊天系统真正难的地方往往不在某一个页面或者某一个接口本身，而在于多条链路同时工作时能不能保持一致。私聊和聊天室都不是前端把消息回显出来就算结束，还要继续确认消息有没有写入数据库、Socket 有没有及时推送、刷新页面后历史记录能不能重新读回来。桌面端提醒也是同样的问题，弹出通知只是最表面的结果，后面还牵涉未读数、托盘状态和窗口焦点变化。如果只看单独一个点，很多问题都不算复杂；放到完整系统里之后，前后端联动和状态同步的影响就会一起出现。

另一个比较深的体会来自 AI 模块。开始做之前，很容易把它理解成“接入一个模型接口”，真正落到系统里才发现关键并不只在模型返回内容本身，而在于历史消息有没有进入索引、检索结果能不能贴近当前讨论、AI 回复生成后能不能继续以普通消息的形式保存和回看。聊天室里的 \texttt{@AI} 不是一个孤立入口，它必须和现有消息流共存，这一点和对话历史管理型问答系统强调的上下文连续性是一致的\cite{historyragqa2026}。这一部分做完以后，我对 AI 接入业务系统的认识比之前具体得多，也更清楚模型能力、消息数据和场景设计之间是如何互相影响的。

从个人收获来看，这次毕设最大的价值在于把一个完整项目真正从头做到尾。以往课程作业通常围绕单个功能展开，范围比较有限；这次需要同时考虑需求分析、系统设计、代码实现、测试验证和结果说明。做完整个过程后，我对接口设计、实时通信、数据库建模、前端状态管理以及问题定位都有了更直观的认识，这些经验比单独完成某一页界面更有价值。

\section{系统不足与展望}

当前版本已经能够支撑课程项目演示和小规模使用，但距离长期稳定运行的技术交流平台还有明显差距。首先，系统目前主要在本地和单机环境下开发与测试，更适合原型验证。如果后续用户数量继续增加，服务端在并发连接、消息广播和 AI 请求调度上的压力会更明显。下一步如果继续扩展，部署方式、缓存机制和服务拆分都需要重新考虑。

其次，AI 模块虽然已经跑通了问答和总结的基本链路，但它对历史语料数量和检索结果质量依赖较强。当聊天室消息较少，或者检索结果不够准确时，回答仍然会偏通用；长上下文利用效率不稳定时，这个问题还会更明显\cite{longcontext2024}。后续更值得投入的方向，不是单纯再换一个模型，而是继续优化消息索引、检索排序和上下文拼接策略，让 AI 返回内容更贴近房间里的真实讨论\cite{ragllmsurvey2024}。即便模型能力继续扩展，大语言模型在业务系统中的稳定性和可控性仍然需要长期处理\cite{llmsurvey2023}。

再次，系统在稳定性方面还有需要补强的地方。第六章测试已经发现，收藏模块在组合筛选条件下仍会出现 500 错误。这说明主路径虽然已经跑通，但复杂路径下的异常处理还不够扎实。相比继续叠加新功能，后续更应该优先补上异常处理、回归测试和日志分析，把现有链路先做稳。

最后，从产品方向看，当前系统已经覆盖了聊天、收藏和 AI 辅助三条主线，但对技术团队更深层次的协作支持还不够。例如，多人代码协同、任务跟踪、知识归档的结构化整理等能力，目前还没有形成完整方案。后续如果继续推进，可以在保持轻量化的前提下逐步补充这些能力，使系统从“技术讨论工具”进一步发展为“面向技术交流的协作平台”。

总体来看，Mini-Chat-Bar 已经验证了一个比较清楚的方向：技术交流系统不应只负责消息传递，还应支持历史回看、结果沉淀和适度的 AI 辅助。当前版本还算不上成熟产品，但主框架已经建立起来，后续只要继续围绕稳定性、检索质量和协作深度推进，系统仍然有较大的完善空间。
